\chapter{Patterns: the architecture in practice}
\label{ch:patterns}

\begin{aside}
The architecture chapters showed primitives. Real apps combine
the primitives in patterns that recur across many projects. This
chapter is a tour: ten common patterns, each with the
architectural shape, the trade-offs, and the code skeleton.
\end{aside}

\section{Pattern 1: master-detail}

\textbf{Shape.} A list on the left, a detail view on the right.
Selecting an item updates the detail.

\textbf{Model.}

\begin{lstlisting}
struct Model {
    std::vector<Item> items;
    int selected = 0;
};
\end{lstlisting}

\textbf{View.}

\begin{lstlisting}
Element view(const Model& m) {
    auto list = render_list(m.items, m.selected);
    auto detail = m.items.empty()
        ? text(\"(no items)\") | Dim
        : render_detail(m.items[m.selected]);
    return h(list | width<30>, detail | grow<1>);
}
\end{lstlisting}

\textbf{Update.} Arrow keys move \code{selected}; Enter triggers
a detail-specific action.

\textbf{Trade-offs.}

\begin{itemize}[leftmargin=*]
  \item Simple to write; works for small datasets.
  \item Re-renders the detail on every selection change; OK as
    long as detail rendering is cheap.
\end{itemize}

\section{Pattern 2: tabbed pages}

\textbf{Shape.} Several distinct pages, switched by Tab or by
keystrokes. Each page has its own state.

\textbf{Model.}

\begin{lstlisting}
struct Model {
    int active_tab = 0;
    HomePage home;
    InboxPage inbox;
    SettingsPage settings;
};
\end{lstlisting}

\textbf{View.}

\begin{lstlisting}
Element view(const Model& m) {
    return v(
        render_tab_bar(m.active_tab),
        switch_tab(m.active_tab,
            home_view(m.home),
            inbox_view(m.inbox),
            settings_view(m.settings))
    );
}
\end{lstlisting}

\textbf{Trade-offs.}

\begin{itemize}[leftmargin=*]
  \item Each page keeps state across tab switches.
  \item Memory cost: all pages' models live even when not
    visible.
  \item Pages can be tested independently.
\end{itemize}

\section{Pattern 3: wizard / stepper}

\textbf{Shape.} Several screens in sequence; each collects input,
the final screen submits.

\textbf{Model.}

\begin{lstlisting}
enum class Step { Welcome, Name, Email, Confirm, Done };
struct Model {
    Step step = Step::Welcome;
    std::string name, email;
};
\end{lstlisting}

\textbf{Update.} Next and Back transitions step the enum;
typing fills fields.

\textbf{Trade-offs.}

\begin{itemize}[leftmargin=*]
  \item Clean linear flow; easy to test (one path per step).
  \item Adding a step is local: enum value, view branch, update
    branch.
\end{itemize}

\section{Pattern 4: undo/redo}

\textbf{Shape.} Every state-changing message creates a new model
that is pushed onto an undo stack.

\textbf{Model.}

\begin{lstlisting}
struct Doc { /* the actual document state */ };
struct Model {
    Doc current;
    std::vector<Doc> undo_stack;
    std::vector<Doc> redo_stack;
};
\end{lstlisting}

\textbf{Update.} Each mutation pushes the previous \code{current}
to \code{undo\_stack}; an \code{Undo} message pops back, pushing
the now-current state to \code{redo\_stack}.

\textbf{Trade-offs.}

\begin{itemize}[leftmargin=*]
  \item Memory grows with the undo history. Cap the stack length
    to bound memory.
  \item Trivially correct because models are values --- no
    fiddly ``do I have a deep copy?''.
\end{itemize}

\section{Pattern 5: incremental search}

\textbf{Shape.} A search box and a filtered result list; results
update on every keystroke without blocking.

\textbf{Model.}

\begin{lstlisting}
struct Model {
    std::string query;
    std::vector<Item> all_items;
    std::vector<size_t> filtered_indices;
    int debounce_token = 0;
};
\end{lstlisting}

\textbf{Update.} A keystroke updates \code{query} and issues
\code{Cmd::after(50ms, Filter\{debounce++\})}. The filter, when
it fires, checks whether \code{debounce\_token} matches; if not,
the keystroke fired again and the filter is stale.

\textbf{Trade-offs.}

\begin{itemize}[leftmargin=*]
  \item Debouncing avoids re-filtering on every key.
  \item Stale-result rejection avoids race-y UI flashes.
  \item More complex \code{update} logic; tests are essential.
\end{itemize}

\section{Pattern 6: side-effecting buttons}

\textbf{Shape.} A button that, when clicked, performs an effect
(network call, file write) and updates UI based on the result.

\textbf{Model.}

\begin{lstlisting}
enum class Status { Idle, Pending, Ok, Failed };
struct Model { Status status = Status::Idle; std::string error; };
\end{lstlisting}

\textbf{Update.} On click, set status to \code{Pending}, issue
\code{Cmd::task}. On result, transition to \code{Ok} or
\code{Failed}.

\textbf{View.} The button shows ``Save'', ``Saving\ldots'',
``Saved!'', or ``Failed (retry)'' depending on status.

\textbf{Trade-offs.}

\begin{itemize}[leftmargin=*]
  \item User sees feedback immediately.
  \item No blocking on the main thread.
  \item Status enum makes the four states explicit; no
    \code{is\_loading} + \code{is\_error} bool combinations.
\end{itemize}

\section{Pattern 7: long-running stream}

\textbf{Shape.} A widget consumes a stream of events (log lines,
tokens from an LLM, network packets) and appends them to a
display.

\textbf{Model.}

\begin{lstlisting}
struct Model {
    std::vector<std::string> lines;
    ScrollState scroll;
    bool autoscroll = true;
};
\end{lstlisting}

\textbf{Update.} Each incoming line via \code{Cmd::task}
(reading from a pipe) is appended. If \code{autoscroll}, adjust
the offset to follow.

\textbf{Trade-offs.}

\begin{itemize}[leftmargin=*]
  \item Use virtualisation (Chapter~\ref{ch:scroll}) for
    large logs.
  \item Cap the vector length to bound memory.
  \item Autoscroll-while-following is a UX pattern users love;
    pair with ``stop following on manual scroll''.
\end{itemize}

\section{Pattern 8: keyboard-driven palette}

\textbf{Shape.} A command palette (Ctrl-P style): a text field, a
filtered list, Enter to execute the selected command.

\textbf{Model.}

\begin{lstlisting}
struct Command { std::string label; std::function<Cmd<Msg>()> action; };
struct Model {
    bool open = false;
    std::string query;
    int selected = 0;
    std::vector<Command> commands;
};
\end{lstlisting}

\textbf{Update.} Open/close on Ctrl-P; arrow keys move
selection; Enter dispatches the command's action.

\textbf{Trade-offs.}

\begin{itemize}[leftmargin=*]
  \item Commands are first-class values; new ones add as a
    push\_back.
  \item Filter is a simple substring match; full-text fuzzy
    match is a deeper rabbit hole.
\end{itemize}

\section{Pattern 9: split-pane editor}

\textbf{Shape.} Two editors side by side; focus moves between
them with a hotkey.

\textbf{Model.}

\begin{lstlisting}
struct Model {
    textarea::Model left;
    textarea::Model right;
    bool focus_right = false;
};
\end{lstlisting}

\textbf{Update.} Keystrokes route to the focused editor; Tab
toggles focus.

\textbf{Trade-offs.}

\begin{itemize}[leftmargin=*]
  \item Two editors compose by message-tagging
    (\code{ToLeft}/\code{ToRight}).
  \item Resizing splits the width evenly; flex handles it.
\end{itemize}

\section{Pattern 10: themed multi-screen app}

\textbf{Shape.} A full app combining several of the patterns
above: tabs at the top, a sidebar, a main area with master-
detail, a status bar, a palette overlay.

\textbf{Composition.}

\begin{lstlisting}
struct App {
    struct Model {
        Theme theme;
        int active_tab;
        SidebarModel sidebar;
        MasterDetailModel main;
        StatusModel status;
        std::optional<PaletteModel> palette;
    };
    // ... Msg with one alt per subsystem
    // ... update routes by alt
    // ... view composes the layout
    // ... subscribe collects every subsystem's subs
};
\end{lstlisting}

\textbf{Trade-offs.} The model is a struct of structs; the message
is a variant of variants; the update is a visitor that routes
to subsystems. The complexity is real but bounded: each
subsystem stays small and testable, and the wiring is
mechanical.

\section{When patterns conflict}

Two patterns might pull in opposite directions. Examples:

\begin{itemize}[leftmargin=*]
  \item \textbf{Undo + long-running stream.} A growing log
    polluting the undo stack is a memory leak. Skip
    appended-log messages from the undo recording.
  \item \textbf{Tabs + palette.} The palette is a modal; while
    open, tab switching is disabled. Handle modality at the
    top-level update.
  \item \textbf{Wizard + autosave.} A wizard step that goes
    back must restore the prior step's input. Save on every
    transition.
\end{itemize}

Patterns are tools, not rules. Read the situation; pick the
right combination.

\section{Anti-patterns}

\begin{warning}
\textbf{God model.} A single \code{Model} with hundreds of
fields. Symptom: \code{update} grows to thousands of lines.
Refactor by splitting into sub-models, one per concern.
\end{warning}

\begin{warning}
\textbf{Message explosion.} A \code{Msg} variant with hundreds
of alternatives. Symptom: every keystroke routes through a
massive visitor. Use sub-component patterns to push detail into
sub-modules.
\end{warning}

\begin{warning}
\textbf{Imperative inside view.} \code{view} calls a function
that performs file I/O. Symptom: the program hangs on render.
Move the I/O into a \code{Cmd::task}.
\end{warning}

\begin{warning}
\textbf{Subscribe of doom.} A monolithic \code{subscribe} with
dozens of always-on subscriptions. Symptom: removing one event
source causes mysterious freezes. Make subscriptions conditional
on relevant model state.
\end{warning}

\section*{Workshop}

\begin{enumerate}[leftmargin=*]
  \item Pick one of the ten patterns and build a small app
    around it.
  \item Combine two patterns (tabs + master-detail) in one
    app. Note where the patterns share state and where they
    are independent.
  \item Take a god-model from an old codebase and refactor it
    into three sub-models. The new \code{update} routes by
    message variant.
\end{enumerate}

\begin{recap}
Patterns are recurring combinations of the architectural
primitives. Master-detail, tabbed pages, wizards, undo, search,
side-effecting buttons, log streams, palettes, split panes,
multi-screen apps --- each combines a model shape, an update
strategy, and a view. Pick patterns intentionally; combine them
carefully; avoid the anti-patterns.
\end{recap}
